% =====================================================================
\section{Background crash course}\label{sec:background}
% =====================================================================

This section teaches, from zero, every idea the project uses. Each subsection ends with where the idea lives
in the code.

% ---------------------------------------------------------------------
\subsection{Fields and partial derivatives}

A \textbf{field} is a quantity that has a value at every place and time, written $u(x,t)$: the temperature
along a rod, the concentration of dye in a pipe, the density of a population along a coastline. Here $x$ is
one-dimensional (a line), so the whole field can be drawn as a heat-map with $x$ on one axis and $t$ on the
other.

A \textbf{partial derivative} measures how $u$ changes when you nudge \emph{one} variable and freeze the other.
The project writes them with subscripts, exactly as the code does:

\begin{center}
\begin{tabular}{@{}llp{9.2cm}@{}}
\toprule
\textbf{Code name} & \textbf{Maths} & \textbf{Meaning} \\ \midrule
\fn{u_t}   & $\partial u/\partial t$     & how fast the value at a fixed point is changing in time \\
\fn{u_x}   & $\partial u/\partial x$     & the slope of the profile at a fixed time (which way is ``uphill'') \\
\fn{u_xx}  & $\partial^2 u/\partial x^2$ & the curvature: positive in a valley, negative on a peak \\
\fn{u_xxx} & $\partial^3 u/\partial x^3$ & third derivative (in the library, never selected) \\
\bottomrule
\end{tabular}
\end{center}

\begin{intuition}[: curvature drives diffusion]
On a peak, the neighbours are lower than the point, so the curvature $u_{xx}$ is negative. In a valley the
neighbours are higher and $u_{xx}$ is positive. ``Move towards the average of your neighbours'' therefore
means ``change in time in proportion to the curvature'': $u_t = D\,u_{xx}$. Peaks sink, valleys fill, the
profile smooths out.
\end{intuition}

% ---------------------------------------------------------------------
\subsection{Partial differential equations (PDEs) and the four mechanism families}

A \textbf{PDE} is a rule connecting a field's partial derivatives. All the equations in this project are
\emph{evolution equations}: they say what $u_t$ is, given the current shape of $u$. The project works with four
mechanism families (and combinations of them):

\begin{center}
\small
\begin{tabular}{@{}p{3.0cm}p{4.2cm}p{7.6cm}@{}}
\toprule
\textbf{Family} & \textbf{Equation} & \textbf{What you would see} \\ \midrule
Diffusion & $u_t = D\,u_{xx}$ & Peaks flatten and spread; sharp features vanish fastest. $D>0$ is the diffusivity. \\
Advection & $u_t + c\,u_x = 0$ & The whole shape slides sideways at speed $c$ without changing shape. \\
Reaction (logistic) & $u_t = r\,u\,(1-u/K)$ & Each point grows (or decays) on its own towards a ceiling $K$; neighbours do not interact. \\
Advection--diffusion & $u_t + c\,u_x = D\,u_{xx}$ & The shape slides \emph{and} spreads: smoke drifting in wind. \\
Reaction--diffusion & $u_t = D\,u_{xx} + r\,u(1-u/K)$ & Spreading plus local growth: a population expanding into new territory. \\
\bottomrule
\end{tabular}
\end{center}

The letters $D, c, r, K$ are \textbf{parameters}: unknown constants that must be estimated from data. A
\textbf{hypothesis} in this project is an equation \emph{shape} (which terms appear) together with its unknown
parameters.

\begin{lesson}[: mechanism = which terms are present]
Identifying the mechanism means identifying \emph{which terms} the equation contains ($u_{xx}$? $u_x$? $u$ and
$u^2$?). Estimating parameters means finding their \emph{numbers}. The project scores these two separately,
because getting the numbers right does not prove you have the right terms (Section~\ref{sec:p7}).
\end{lesson}

% ---------------------------------------------------------------------
\subsection{Initial and boundary conditions, and well-posedness}

An evolution equation alone does not determine the future; you also need:
\begin{itemize}
  \item an \textbf{initial condition} (IC): the profile at $t=0$. Here
        $u(x,0) = \sin(\pi x) + 0.5\sin(3\pi x)$, a ``double hump'';
  \item \textbf{boundary conditions} (BC): what happens at the ends of the rod. Here $u(0,t)=u(1,t)=0$, the ends
        are held at zero (\emph{Dirichlet} conditions).
\end{itemize}
A problem is \textbf{well-posed} if a solution exists, is unique, and changes only a little when the data changes
a little. Two physical facts used later:
\begin{itemize}
  \item Diffusion with $D<0$ is the \emph{backward heat equation}: it un-mixes, and tiny wiggles explode. It is
        ill-posed, so the project enforces the physical bound $D>0$.
  \item Pure advection carries information in one direction only, so it can accept a boundary value only at the
        inflow end. Prescribing values at \emph{both} ends over-determines it; the project warns about this.
\end{itemize}

% ---------------------------------------------------------------------
\subsection{Simulating a PDE: finite differences}

Computers cannot handle a continuum, so the rod is cut into grid points $x_i$ spaced $\Delta x$ apart and time
into steps $\Delta t$. Derivatives become differences of neighbouring values:
\[
  u_{xx}(x_i) \approx \frac{u_{i+1} - 2u_i + u_{i-1}}{\Delta x^2}, \qquad
  u^{n+1}_i = u^n_i + \Delta t\, D\, \frac{u^n_{i+1} - 2u^n_i + u^n_{i-1}}{\Delta x^2}.
\]
The second formula is the \textbf{explicit (forward-time, centred-space) scheme} used to generate the data in
Phase~1. It is only stable if $D\,\Delta t/\Delta x^2 \le 1/2$, which is why the simulation uses 800 time
steps for 61 grid points. Later phases also need a general solver for \emph{any} candidate equation
(\fn{validation/forward_solver.py}); it offers simple Euler steps and the more accurate fourth-order
Runge--Kutta (RK4) method. Euler with centred differences is unconditionally unstable for advection, which is
why RK4 was added in Phase~7.

% ---------------------------------------------------------------------
\subsection{Numerical derivatives from data, and why noise is the enemy}

To \emph{discover} an equation from data you need the derivatives of the data. With scattered points you first
\textbf{reconstruct} a regular grid by interpolation, optionally \textbf{smooth} it, then take finite
differences. The trouble is that differencing amplifies noise. If each value has noise of size $\sigma$, the
second-difference formula divides by $\Delta x^2$, so its error is roughly
\[
  \text{error in } u_{xx} \;\sim\; \frac{\sigma}{\Delta x^2},
\]
which becomes enormous on a fine grid. A third derivative is worse still. Smoothing reduces the noise but
introduces bias. This single fact explains why Phase~2 works on clean data and fails on noisy data, and why the
project later turned to PINNs, which never take numerical differences of the data.

% ---------------------------------------------------------------------
\subsection{Sparse regression and SINDy}\label{sec:sindy}

\textbf{SINDy} (Sparse Identification of Nonlinear Dynamics) turns ``find the equation'' into a regression:
\begin{enumerate}
  \item Compute $u_t$ at many points: this is the target vector $y$.
  \item Build a \textbf{library} $\Theta$ whose columns are candidate terms evaluated at the same points. Here:
        $1,\ u,\ u_x,\ u_{xx},\ u_{xxx},\ u^2,\ u\,u_x,\ u\,u_{xx}$ (8 columns).
  \item Find coefficients $\xi$ with $y \approx \Theta\,\xi$, while forcing most entries of $\xi$ to be exactly
        zero, because real laws have few terms.
\end{enumerate}
The project's sparsifier is \textbf{STLSQ} (Sequentially Thresholded Least Squares): fit by ordinary least
squares, set every coefficient smaller than a threshold (0.05) to zero, refit using only the surviving
columns, and repeat until nothing changes. Columns are normalised first so the threshold means the same thing
for every term.

\begin{caution}[: identifiability and the eigenfunction trap]
Regression can only separate terms whose columns look different. For $\sin(\pi x)$, the second derivative is
$u_{xx} = -\pi^2 u$ \emph{exactly}, so the $u$ and $u_{xx}$ columns are perfectly proportional. No method can
tell $D\,u_{xx}$ apart from $-D\pi^2 u$. Phase~2 fell into exactly this trap with the original Phase~1 data and
``discovered'' $u_t = -0.987\,u$: a perfect fit, with the wrong physics. Using two sine modes with different
curvatures fixed it.
\end{caution}

% ---------------------------------------------------------------------
\subsection{Neural networks, automatic differentiation, and training}

A \textbf{multilayer perceptron} (MLP) is a function built from layers: each layer multiplies its input by a
weight matrix, adds a bias, and applies a squashing function (here $\tanh$). The project's network takes the two
numbers $(x,t)$ and outputs one number $u_\theta(x,t)$, where $\theta$ stands for all the weights. It has four
hidden layers of 64 units.

\textbf{Training} means adjusting $\theta$ to make a loss (an error score) small, using its gradient. The project
uses \textbf{Adam} (3000 steps, a robust first-order optimiser) and then \textbf{L-BFGS} (300 iterations, a
quasi-Newton method that converges much further once Adam is close).

\textbf{Automatic differentiation} (autograd, \fn{torch.autograd.grad}) gives \emph{exact} derivatives of the
network's output with respect to its inputs. So $\partial u_\theta/\partial t$ and
$\partial^2 u_\theta/\partial x^2$ come for free, with no finite differences and no noise amplification. This is
the key to PINNs.

% ---------------------------------------------------------------------
\subsection{Physics-informed neural networks (PINNs)}\label{sec:pinn}

A \textbf{PINN} is a neural network trained to do two things at once: pass close to the observations
\emph{and} obey a candidate equation everywhere. For a candidate written as a residual $\mathcal{R}=0$ (for
diffusion, $\mathcal{R} = u_t - D\,u_{xx}$) the loss is
\[
  \mathcal{L} \;=\;
  \underbrace{\tfrac1N\textstyle\sum_i \big(u_\theta(x_i,t_i) - u_i\big)^2}_{\text{data: fit observations}}
  + \underbrace{\tfrac1M\textstyle\sum_j \mathcal{R}(x_j,t_j)^2}_{\text{physics: obey the equation}}
  + \underbrace{\mathcal{L}_{\text{BC}}}_{\text{ends held at 0}}
  + \underbrace{\mathcal{L}_{\text{IC}}}_{\text{start profile}} .
\]
The physics term is evaluated at random \textbf{collocation points} $(x_j,t_j)$ that need no measurements.
All four weights are 1 in this project.
\begin{itemize}
  \item \textbf{Forward PINN:} the parameters are known (e.g.\ $D=0.1$ fixed). The PINN acts as a smart
        interpolator/denoiser that respects physics.
  \item \textbf{Inverse PINN:} the parameters are unknown. $D$ becomes a trainable number, optimised together
        with $\theta$. Whatever $D$ lets the network fit data \emph{and} physics best is the estimate.
\end{itemize}

\begin{intuition}[: why a PINN can reject a wrong equation]
If a candidate equation is wrong, no network can satisfy both the data and that equation at the same time;
one of the loss terms stays large. If it is right, all terms can be driven to near zero. So the final losses of a
trained inverse PINN are \emph{evidence} for or against the hypothesis, and the fitted parameters are estimates.
\end{intuition}

% ---------------------------------------------------------------------
\subsection{Large language models as hypothesis generators}

A \textbf{large language model} (LLM) such as Google Gemini predicts text. Given a description of what was
observed (``the amplitude decreased substantially, the peak did not move\ldots''), it can propose plausible
equations. The project treats the LLM strictly as a \emph{brainstormer}:
\begin{itemize}
  \item it must answer in \textbf{structured JSON} matching a schema (\fn{Hypothesis}: id, name, equation,
        parameters, required derivatives, mechanism, assumptions, rationale, confidence);
  \item its \fn{confidence} field is recorded but \emph{never} used as evidence;
  \item a deterministic validator checks syntax, variables, parameters and derivatives \emph{before} any
        expensive PINN training;
  \item the verdict comes only from measured PINN losses.
\end{itemize}
For testing without internet access, a \textbf{mock provider} returns a fixed, canned set of answers. All
Phase~4--7 numbers come from the mock; Section~\ref{sec:gemini} reports the real Gemini run.

% ---------------------------------------------------------------------
\subsection{Retrieval-augmented generation (RAG)}

\textbf{RAG} means: before asking the LLM, retrieve relevant documents and include them in the prompt, so its
proposals are grounded in evidence. The steps:
\begin{enumerate}
  \item \textbf{Corpus:} five short Markdown ``textbook'' documents, one per mechanism.
  \item \textbf{Chunking:} split each document at its section headers into chunks.
  \item \textbf{Embedding:} turn each chunk into a vector. The project uses a simple, deterministic
        \emph{bag-of-words hashing} embedding: each word is hashed into one of 256 slots, slot counts form the
        vector, and the vector is normalised to length 1.
  \item \textbf{Search:} embed the query (the observation description) the same way and return the chunks with
        the highest \emph{cosine similarity} (the dot product of unit vectors).
\end{enumerate}
\textbf{Answer leakage} is the danger that retrieved text gives away the answer. The corpus states only generic
textbook laws (``Fick's second law, $u_t = D\,u_{xx}$'') and never this dataset's value $D=0.1$; an automated
check verifies this.

% ---------------------------------------------------------------------
\subsection{Knowledge graphs}

A \textbf{knowledge graph} stores facts as typed nodes and typed edges: \emph{Mechanism} ``diffusion''
\textsc{governs} \emph{Equation} ``$u_t=D\,u_{xx}$'', which \textsc{has\_parameter} \emph{Parameter} ``$D$'',
\textsc{supported\_by} \emph{Document} ``diffusion.md''. Every edge records its \textbf{provenance} (which
document and chunk it came from). Where RAG answers ``what text is relevant?'', the graph answers structured
questions such as ``which parameters appear in the documented equations for diffusion?'' (answer: $D$). In this
project the graph has 22 entities and 40 relationships, all sourced, and it supplies \emph{equation templates}
that join the LLM's proposals in the candidate pool.

% ---------------------------------------------------------------------
\subsection{Choosing between hypotheses: splits, parsimony, BIC}\label{sec:selection}

\textbf{Data splits.} To judge a model honestly you must test it on data it never saw:
\begin{itemize}
  \item \textbf{train:} used to fit (observations with $t \le 0.8$, minus a holdout);
  \item \textbf{validation (held-out):} a random 20\% of the $t\le0.8$ points, used for the accept/reject
        \emph{gate} and for selection;
  \item \textbf{OOD (out-of-distribution):} all points with $t>0.8$, never used to train or select. It
        measures genuine extrapolation into the future.
\end{itemize}

\textbf{Nested models.} Reaction--diffusion with $r=0$ \emph{is} diffusion; advection--diffusion with $c=0$
\emph{is} diffusion. A more general model can always fit at least as well, so ``best fit'' alone favours
needless complexity. This is the \emph{nested-model problem}.

\textbf{Parsimony (Occam's razor).} Among models that fit equally well, prefer the simplest. The project makes
this precise:
\begin{enumerate}
  \item \emph{Gate:} reject a candidate unless its held-out error, physics residual, BC/IC losses and physical
        bounds are all within fixed thresholds.
  \item \emph{Equivalence:} keep the gate survivors whose held-out RMSE is within 5\% of the best.
  \item \emph{Complexity:} choose the lowest score
        $=\#\text{terms} + \#\text{free parameters} + \text{max derivative order} + 2\cdot\#\text{nonlinear terms}
        + 0.1\cdot\#\text{operations}$.
  \item \emph{BIC} as a tie-breaker and cross-check: $\text{BIC} = n\ln(\text{MSE}) + k\ln n$, which rewards
        fit and charges $\ln n$ per parameter.
\end{enumerate}

\textbf{Term necessity.} For each parameter, measure what share of $u_t$ (RMS) its terms contribute on the
trained PINN. Below 5\% the parameter is \emph{inactive}: the data does not need that term. Phase~7 promoted
this into the selection rule (Section~\ref{sec:p7}).

\textbf{The oracle rule.} Because the data is synthetic, the clean truth is known. It may be used \emph{only}
to report error \emph{after} every decision has been made (the function \fn{oracle_evaluation}), never to make a
decision. Phase~7 audits the code automatically for violations of this rule.

\begin{incode}
Derivatives: \fn{equation_discovery/derivatives.py} (finite differences) and \fn{pinn/derivatives.py}
(autograd). SINDy: \fn{equation_discovery/sindy.py}. PINN: \fn{pinn/}. LLM: \fn{llm/}. RAG: \fn{rag/}.
Graph: \fn{knowledge_graph/}. Selection: \fn{validation/model_selection.py}, \fn{hypotheses/complexity.py}.
Splits and the gate: \fn{hypotheses/search.py}.
\end{incode}
